Die Abbildung
\maabbdisp {u} { I } {\R^n
} {}
wird bezüglich der Basischnitte
\mathl{e_1 , \ldots , e_n}{} durch

\mavergleichskettedisp
{\vergleichskette
{ u
}
{ =} { u_1e_1 + \cdots + u_ne_n
}
{ } {
}
{ } {
}
{ } {
}
}
{}{}{}
beschrieben. Die
\definitionsverweis {vertikale Ableitung}{}{}
des Zusammenhangs dieses Schnittes in Richtung der einzigen Ableitungsrichtung $\partial$ ist nach
[[Differenzierbare Mannigfaltigkeit/R/Vektorbündel/Linearer Zusammenhang/Lokale Beschreibung/Fakt|Fakt]]
gleich

\mavergleichskettedisp
{\vergleichskette
{  \sum_{k = 1}^n \left(    \partial u_k  +  \sum_{j = 1}^n u_j \Gamma_{j}^k  \right) e_k
}
{ =} { \sum_{k = 1}^n \left(   u_k'  +  \sum_{j = 1}^n u_j \Gamma_{j}^k  \right) e_k
}
{ } {
}
{ } {
}
{ } {
}
}
{}{}{.}
Dabei ist $u$ nach
[[Mannigfaltigkeit/Vektorbündel/R/Zusammenhang/Horizontaler Schnitt/Vertikale Ableitung/Fakt|Fakt]]
genau dann ein horizontaler Schnitt, wenn die vertikale Ableitung gleich $0$ ist, und dies bedeutet für die einzelnen Komponenten, dass

\mavergleichskettedisp
{\vergleichskette
{ u_k' + \sum_{j = 1}^n \Gamma_j^k u_j
}
{ =} { 0
}
{ } {
}
{ } {
}
{ } {
}
}
{}{}{}
für alle $k$ gelten muss. Dies ist äquivalent zu

\mavergleichskettedisp
{\vergleichskette
{ u_k'
}
{ =} { \sum_{j = 1}^n a_{kj} u_j
}
{ } {
}
{ } {
}
{ } {
}
}
{}{}{}
für alle $k$, was gerade bedeutet, dass $u$ das Matrixsystem

\mavergleichskettedisp
{\vergleichskette
{ u'
}
{ =} { \left(  a_{kj}    \right) u
}
{ } {
}
{ } {
}
{ } {
}
}
{}{}{}
erfüllt, also eine Lösung des linearen Differentialgleichungssystems ist.

 [[Kategorie:Latexseite]]
